\section[WHAT THE RECORD IDENTIFIES, AND WHAT IT ONLY BOUNDS]{WHAT THE RECORD IDENTIFIES,\\ AND WHAT IT ONLY BOUNDS}\label{sec:bounds}

The finding record does not point-identify the underlying condition, but it
does impose useful bounds. Three statements hold, needing in turn no
additional assumption about detection, one assumption about the direction of
follow-up, and one number about its size. All three rest on the
no-false-positive assumption of Section~\ref{sec:intro}.

\subsection*{With no additional assumption about detection: a floor and a wide range}

Because detection cannot exceed 100 percent, the true problem rate in each
history group must be at least as high as the recorded finding rate: $d_f \le
1$ in equation~\eqref{eq:qrd} gives $r_f \ge q_f$. If five percent of
clean-history program-years record a finding, at least five percent had the
condition. Without known complete detection, a recorded finding rate is a lower bound
on the underlying condition rate, not an estimate of it.

The question a planner asks is about the ratio, and the floor alone leaves it
wide. $\Rlat$ is largest when the prior-finding group's problem rate is as
high as possible (one) and the clean group's is as low as the record allows
($q_0$), and smallest in the reverse case, so
%
\begin{equation}\label{eq:unres}
  \Rlat \;\in\; [\,q_1,\; 1/q_0\,] ,
\end{equation}
%
and both ends can actually occur (Proposition~\ref{prop:bounds}). At the
values in Table~\ref{tab:worlds}, $q_0 = \nWorldQzero{}$ and $q_1 =
\nWorldQone{}$, so $\Rlat$ can range from $\nWorldQone{}$ to $20$. The same
fourfold recorded ratio is consistent with the underlying condition being five
times less common, equally common, or much more common after a prior finding.
Without information about detection, the record does not even fix the sign of
the underlying association.

\subsection*{With one assumption: the observed ratio is a ceiling}

Now assume that follow-up never makes the auditor look \emph{less} hard at a
program that had a finding than at one that did not: an existing problem is at
least as likely to be detected after a prior finding as after a clean history,
%
\[
  d_1 \;\ge\; d_0 , \qquad \text{equivalently } M \ge 1 .
\]
%
Call this monotone detection. The direction is consistent with a follow-up
regime, which directs more attention toward prior findings and never less,
but monotone detection remains a substantive assumption about how follow-up
affects the probability of detecting an existing condition, not a fact the
regulation proves; it fails, for example, when a prior finding shifts
attention toward one area and away from another, and it may hold for some
kinds of condition and not others.

Under monotone detection, dividing equation~\eqref{eq:product} by $M \ge 1$
gives
%
\begin{equation}\label{eq:mono}
  \Rlat \;\le\; \Robs .
\end{equation}
%
If follow-up can only increase detection, the recorded risk ratio cannot be
smaller than the underlying risk ratio. In Table~\ref{tab:worlds}, the underlying ratio is at most 4 in every
column, and the columns show it taking the values 4, 1 and 0.5. One
consequence needs no further information: if the recorded ratio is below one,
so that programs with a prior finding record \emph{fewer} findings than
programs without, the underlying association is genuinely negative, and no
amount of differential follow-up can explain it away.

\subsection*{With one number: a floor on the underlying ratio}

The constructive result comes from putting a ceiling on the follow-up
detection multiplier. Write $\Mbar$ for the largest value of $M$ the
organization is prepared to defend: the most that follow-up could plausibly
raise the chance of detecting an existing problem. If $1 \le M \le \Mbar$,
dividing equation~\eqref{eq:product} by $M \le \Mbar$ gives
%
\begin{equation}\label{eq:mbar}
  \Rlat \;\ge\; \Robs / \Mbar .
\end{equation}
%
Even after giving follow-up the largest detection advantage the organization
considers plausible, the underlying condition is still at least $\Robs/\Mbar$
times as common after a prior finding. A recorded ratio of 4 and a defensible
belief that follow-up can increase detection by no more than a factor of 2
implies an underlying risk ratio of at least 2. With equation~\eqref{eq:mono},
the record then places $\Rlat$ between 2 and 4: the inference survives, cut
down by the part of the elevation the follow-up can account for. And whenever
the recorded ratio exceeds $\Mbar$, the underlying association is genuinely
positive (Proposition~\ref{prop:mono}).

These bounds are sharp once the floor $r_f \ge q_f$ is added. Each detection
probability lies between $q_f$ and one, so the multiplier cannot exceed
$1/q_0$, and the multipliers consistent with the record are $1 \le M \le
\min\{\Mbar, 1/q_0\}$. The identified set is
%
\begin{equation}\label{eq:sharp}
  \Rlat \;\in\; \Bigl[\, \max\Bigl\{ q_1,\; \frac{\Robs}{\Mbar} \Bigr\},\; \Robs \,\Bigr] ,
\end{equation}
%
and every value in it is produced by some admissible pair of detection
probabilities (Proposition~\ref{prop:mono}(iii)). At the values of
Table~\ref{tab:worlds} with $\Mbar = 2$ the set is $[2, 4]$.

The binding identification constraint on what an auditor can learn from a
finding record is therefore not sample size or a more elaborate estimator; it
is what can be said about detection, summarized in $M$. Sample size still
governs precision. The propositions are population statements about $q_0$ and
$q_1$; in data, $\Robs$ is estimated, and each conclusion needs the matching
confidence limit. A negative underlying association requires an upper
confidence limit for $\Robs$ below one. A positive association under a
defended $\Mbar$ requires a lower confidence limit for $\Robs$ above
$\Mbar$. What can be said about detection is a question about audit
procedures, answerable by the organization that performs them and not from
the record; Section~\ref{sec:operational} sets out how.

\subsection*{More years of history do not remove the problem}

A longer record does not solve this on its own. Take any recorded history of
any length and any story about how likely the underlying problem was at each
point, so long as it stays at or above the recorded rate. Setting detection
at each point to whatever value reconciles the two reproduces the entire
recorded history exactly (Proposition~\ref{prop:full}). Adding years adds observations of the product
$q = rd$, not of either factor. Restricted models of the underlying condition
can nevertheless be identifiable when substantive restrictions constrain how
the condition evolves and how detection responds to the record
\citep{allman2009}; identification then comes from those restrictions, not
from record length alone. This paper does not estimate such a model. Its
interest is what can be said before those restrictions are imposed, and what
audit design can add.
